\documentclass[10pt,letterpaper]{article}
\usepackage{cornell-notes}

\title{Clause 26.06.02.02: Class Template Empty View}
\author{}
\date{}
\setDocTitle{Clause 26.06.02.02: Class Template Empty View}
\setDocSubtitle{C++ 2024}
\setDocOwner{C++ 2024 Cornell Notes}
\setDocDate{\today}
\setCornellCollection{C++ 2024 Cornell Notes}
\setCornellUnitType{Clause}
\setCornellUnitNumber{26.06.02.02}
\setCornellUnitTitle{Class Template Empty View}
\hypersetup{pdftitle={Clause 26.06.02.02: Class Template Empty View},pdfsubject={Cornell notes},pdfauthor={}}

\begin{document}
\maketitle
\begin{CornellOverviewBox}{Overview}
\textbf{Classification:} Standard library - Normative\\[2pt]
\textbf{Learning objective:} Understand the rules, terminology, and semantic consequences associated with Class template empty\_view.
\end{CornellOverviewBox}\begin{CornellSummaryBox}{Summary}
{\sffamily\bfseries\color{Accent} Summary}\par\vspace{3pt}Section 26.6.2.2 focuses on Class template empty\_view. Use the listed grammar productions together with the semantic constraints and disambiguation rules referenced by the section. When applying it, preserve the distinction between mandatory requirements, recommendations, permissions, and behavior deliberately left undefined, unspecified, or implementation-defined.
\end{CornellSummaryBox}

\end{document}
